% Section content template for individual Educative lessons
% This file contains the actual content for a specific section
% Generated from Educative API section content

% Section: Code for LinkedIn
% Section ID: 6071551264030720
% Section Slug: code-for-linkedin
% Generated: 2025-12-14T13:35:27.461238

We've reviewed LinkedIn's various aspects and observed the attributes attached to the problem using various UML diagrams. Now , let's explore the more practical side of things. We will implement the LinkedIn network using multiple languages. This is usually the last step in an object-oriented design interview process.
\\
We have chosen the following languages to write the skeleton code of the different classes present in LinkedIn:

\begin{itemize}
\item
\protect\phantomsection\label{6FaaQONUvywtbdyHyuXv_}
Java
\item
\protect\phantomsection\label{4rYNo4NzLydhJ1Wnia9DC}
C\#
\item
\protect\phantomsection\label{lqaYRl4VFdv6GsrS8HMHB}
Python
\item
\protect\phantomsection\label{RTGTXRliywNc6SfDsfSe4}
C++
\item
\protect\phantomsection\label{n6NNkz2IYP7UMSX9qB7Yd}
JavaScript
\end{itemize}
\\
If you want to skip directly to the implementation and see the complete workflow, simply scroll down or click~\hyperref[Executable-code-Designing-LinkedIn]{Executable Code: Designing LinkedIn}.

\subsection{LinkedIn classes}\label{HWWWO9BJuj0MQDkqv-iGS}
\\
This section will provide the skeleton code of the classes designed in the class diagram lesson.
\begin{quote}
\textbf{Note:} For simplicity, we are not defining getter and setter functions. The reader can assume that all class attributes are private, accessed through their respective public getter methods, and modified only through their public method functions.
\end{quote}
\subsubsection{Constants}\label{07f6fUZCabYvPLkDPg5J_}
\\
The following code defines the various enums and custom data types being used in the LinkedIn design:
\begin{quote}
\textbf{Note:} JavaScript does not support enumerations, so we will use the \texttt{Object.freeze()} method as an alternative that freezes an object and prevents further modifications.
\end{quote}

\textbf{Constant definitions}

\begin{lstlisting}[language=Java]
public class Address {
  private int zipCode;
  private String streetAddress;
  private String city;
  private String state;
  private String country;
}

enum ConnectionInviteStatus {
  PENDING, 
  ACCEPTED,
  IGNORED
}

enum JobStatus {
  OPEN,
  ON_HOLD,
  CLOSED
}
\end{lstlisting}

\subsubsection{Person and user}\label{KXhfpzME__dgiH9DoEzPZ}
\\
The \texttt{Person} class is an abstract class that contains details such as name, address, phone number, email, and password. It is derived from the \texttt{User} class.

\textbf{The Person, Admin, and User classes}

\begin{lstlisting}[language=Java]
// Person will be an abstract class
public abstract class Person {
  private String name;
  private Address address;
  private String phone;
  private String email;
  private String password;
}

public class User extends Person {
  private int userId;
  private Date dateOfJoining;
  private Profile profile;
  private List<User> connections;
  private List<User> followsUsers;
  private List<CompanyPage> followCompanies;
  private List<Group> joinedGroups;
  private List<CompanyPage> createdPages;
  private List<Group> createdGroups;

  public boolean search(User user);
  public boolean sendMessage(Message message);
  public boolean sendInvite(ConnectionInvitation invite);
  public boolean cancelInvite(ConnectionInvitation invite);
  public Page createPage(CompanyPage page);
  public boolean deletePage(CompanyPage page);
  public Group createGroup(Group group);
  public boolean deleteGroup(Group group);
  public Post createPost(Post post);
  public boolean deletePost(Post post);
  public Comment createComment(Comment comment);
  public boolean deleteComment(Comment comment);
  public boolean react(Post post);
  public boolean requestRecommendation(User user);
  public boolean acceptRecommendation(User user);
  public boolean applyForJob(Job job);
}
\end{lstlisting}

\subsubsection{Recommendation, achievement, and analytics}\label{WZiZyEhPKAWac9OTeO0qN}
\\
The \texttt{Recommendation}, \texttt{Achievement}, and \texttt{Analytics} classes will provide a user's personal information and make up the \texttt{Profile} class. The definition of these classes is given below:

\textbf{The Recommendation, Achievement, and Analytics classes}

\begin{lstlisting}[language=Java]
public class Recommendation {
  private int userId;
  private Date createdOn;
  private String description;
  private boolean isAccepted;
}

public class Achievement {
  private String title;
  private Date dateAwarded;
  private String description;
}

public class Analytics {
  private int searchAppearances;
  private int profileViews;
  private int postImpressions;
  private int totalConnections;
}
\end{lstlisting}

\subsubsection{Profile, experience, education, and skill}\label{nPZwCloSpO2pu8LIKM2W1}
\\
The \texttt{Experience}, \texttt{Education}, and \texttt{Skill} classes will provide a user's personal information and make up the \texttt{Profile} class. The definition of these classes is given below:

\textbf{The Experience, Education, Skill, and Profile classes}

\begin{lstlisting}[language=Java]
public class Experience {
  private String title;
  private String company;
  private Address location;
  private Date startDate;
  private Date endDate;
  private String description;
}

public class Education {
  private String school;
  private String degree;
  private Date startDate;
  private Date endDate;
  private String description;
}

public class Skill {
  private string name;
}

public class Profile {
  private String headline;
  private String about;
  private String gender;
  private List<byte> profilePicture;
  private List<byte> coverPhoto;

  private List<Experience> experiences;
  private List<Education> educations;
  private List<Skill> skills;
  private List<Achievement> achievements;
  private List<Recommendation> recommendations;
  private Analytics analytics;

  public boolean addExperience(Experience experience);
  public boolean addEducation(Education education);
  public boolean addSkill(Skill skill);
  public boolean addAchievement(Achievement achievement);
}
\end{lstlisting}

\subsubsection{Company, job, and group}\label{VRwvdmI5hHTkK2vx5I9yn}
\\
LinkedIn users can create groups and company pages. The company page contains information about the company. The company pages will host various job postings. The
\texttt{Job}, \texttt{CompanyPage}, and \texttt{Group} classes are shown below:

\textbf{The Job, CompanyPage, and Group classes}

\begin{lstlisting}[language=Java]
public class Job {
  private int jobId;
  private String jobTitle;
  private Date dateOfPosting;
  private String description;
  private CompanyPage company;
  private String employmentType;
  private Address location;
  private JobStatus status;
}

public class CompanyPage {
  private int pageId;
  private String name;
  private String description;
  private String type;
  private int companySize;
  private User createdBy;
  private List<Job> jobs;

  public boolean createJobPosting();
  public boolean deleteJobPosting(Job job);
}

public class Group {
  private int groupId;
  private String name;
  private String description;
  private int totalMembers;
  private List<User> members;

  public boolean updateDescription();
}
\end{lstlisting}

\subsubsection{Post, comment, message, and connection invitation}\label{dlqUBwjKZaL-y2_tZ0owV}
\\
LinkedIn users can create posts and comments. They can also send messages and connection invitations to other users. The definition of \texttt{Post}, \texttt{Comment}, \texttt{Message}, and \texttt{ConnectionInvitation} classes is given below:

\textbf{The Post, Comment, Message, and ConnectionInvitation classes}

\begin{lstlisting}[language=Java]
public class Post {
  private int postId;
  private User postOwner;
  private String text;
  private List<byte> media;
  private int totalReacts;
  private int totalShares;
  private List<Comment> comments;

  public boolean updateText();
}

public class Comment {
  private int commentId;
  private User commentOwner;
  private String text;
  private int totalReacts;
  private List<Comment> comments;

  public boolean updateText();
}

public class Message {
  private int messageId;
  private User sender;
  private List<User> recipients;
  private String text;
  private List<byte> media;

  public boolean addRecipients(List<User> recipients);
}

public class ConnectionInvitation {
  private User sender;
  private User recipients;
  private Date dateCreated;
  private ConnectionInviteStatus status;

  public boolean acceptConnection();
  public boolean rejectConnection();
}
\end{lstlisting}

\subsubsection{Search and notification}\label{zERZ8ATAcYeau0QMXGV-8}
\\
The \texttt{Search} class contains information on users, company pages, groups, and jobs. It enables the search functionality based on the given criteria (user, company page, group, and job keywords).
\\
The \texttt{Notification} class is responsible for sending notifications to users about any new messages, comments, posts, or connection invitations via the built-in notification option.
\\
The definition of these classes is given below:

\textbf{The Search interface and Notification classes}

\begin{lstlisting}[language=Java]
public Search {

  public List<User> searchUser(String name) {
    // functionality
  }

  public List<CompanyPage> searchCompany(String name) {
    // functionality
  }

  public List<Job> searchJobs(String title) {
    // functionality
  }

  public List<Group> searchGroups(String name) {
    // functionality
  }
}

public class Notification {
  private int notificationId;
  private Date createdOn;
  private String content;

  public boolean sendNotification(Account account);
}
\end{lstlisting}

\subsection{Executable code: Designing LinkedIn}\label{xs-4c2nQg8LDWvadptkwn}
\\
Below is a fully self-contained, runnable program in Java, C\#, C++, Python, and JavaScript that demonstrates the core workflows of a LinkedIn system. The main driver code resides in \texttt{Driver.java}, \texttt{Driver.cs}, \texttt{Driver.py}, \texttt{Driver.cpp}, and \texttt{Driver.js} for the respective languages. You can click the ``Run'' button to execute the code.


\subsubsection{What does this code show}\label{NvONzLe2yVZ6pBBkuXgUF}

\begin{itemize}
\item
\protect\phantomsection\label{QNhNF0lWcoHL4LCsans-7}
\textbf{System Initialization:} The simulation begins by setting up key components of the LinkedIn system: two users (Alice and Bob) with their respective profiles and addresses. Supporting objects, such as \texttt{Search}, \texttt{ConnectionInvitation}, and \texttt{Notification} are also ready. This sets the stage for simulating common LinkedIn interactions such as searching for users, sending connection invites, and engaging with posts.
\item
\protect\phantomsection\label{Ja9Pb1jptavE7CAZAYazm}
\textbf{Scenario 1 (User searches and sends a connection invitation):} User A (Alice) searches for User B (Bob) using the search feature. When Bob is found, a notification is automatically sent to him indicating he appeared in a search. Alice then creates a connection invitation and sends it to Bob. Another notification is triggered to inform Bob about the new connection request.
\item
\protect\phantomsection\label{37Nusdx2j7vxmAjmHorJJ}
\textbf{Scenario 2 (Connection request is accepted):} User B (Bob) receives the connection invitation from Alice. He chooses to accept the request, which updates the connection status to ``accepted.'' The system then sends a notification to Alice, confirming that her connection request was approved.
\item
\protect\phantomsection\label{Ta_ogv30GFWYKXifp8kp2}
\textbf{Scenario 3 (User interaction via post and comment):} User A (Alice) creates a new post announcing a professional update. User B (Bob) views the post and comments on it. A notification is automatically sent to Alice, informing her that Bob commented on her post.
\end{itemize}
\begin{quote}
\textbf{Hint: Try customizing the code!}
\\
This code is designed not just for reading, but for experimentation. You can~edit, extend, or modify~any part of the example.
\end{quote}

\textbf{Designing LinkedIn}

\begin{lstlisting}[language=Java]
import java.util.*;

public class Driver {

    public static void main(String[] args) {
        // Create common address
        Address address = new Address(12345, "123 Main St", "Cityville", "StateX", "CountryY");

        // Create user profiles
        Profile profileA = new Profile("Software Engineer", "Passionate about tech", "Male");
        Profile profileB = new Profile("Data Analyst", "Loves data", "Female");

        // Create users
        User userA = new User("Alice", address, "1234567890", "alice@example.com", "password", 1, new Date(), profileA);
        User userB = new User("Bob", address, "0987654321", "bob@example.com", "password", 2, new Date(), profileB);

        System.out.println("🔹 Users created:");
        System.out.println("User A: " + userA.getName());
        System.out.println("User B: " + userB.getName());

        // ========================
        // Scenario 1: Search & Connection Request
        // ========================
        System.out.println("
🔸 Scenario 1: User A searches for User B and sends a connection request");

        Search search = new Search();
        List<User> foundUsers = search.searchUser("Bob");

        // Simulate search result
        foundUsers.add(userB);

        if (!foundUsers.isEmpty()) {
            System.out.println("✔ User B found in search results");

            // Auto-notify User B of being searched
            Notification searchNotif = new Notification(1001, "You appeared in a search result.");
            searchNotif.sendNotification(userB);

            // Send connection request
            ConnectionInvitation invite = new ConnectionInvitation(userA, userB);
            userA.sendInvite(invite);
            System.out.println("📨 Connection request sent from " + userA.getName() + " to " + userB.getName());

            // Auto-notify User B of connection request
            Notification inviteNotif = new Notification(1002, userA.getName() + " sent you a connection request.");
            inviteNotif.sendNotification(userB);
        }

        // ========================
        // Scenario 2: User B accepts the connection
        // ========================
        System.out.println("
🔸 Scenario 2: User B accepts the connection request");

        ConnectionInvitation invitation = new ConnectionInvitation(userA, userB);
        boolean accepted = invitation.acceptConnection();

        if (accepted) {
            System.out.println("🤝 " + userB.getName() + " accepted the connection request");
            Notification notifyAccepted = new Notification(1003, userB.getName() + " accepted your connection request.");
            notifyAccepted.sendNotification(userA);
        } else {
            System.out.println("❌ Connection request not accepted");
        }

        // ========================
        // Scenario 3: User A creates a post and User B comments on it
        // ========================
        System.out.println("
🔸 Scenario 3: User A creates a post and User B comments on it");

        Post post = new Post(101, userA, "Excited to start my new job!");
        System.out.println("📝 " + userA.getName() + " posted: "" + post.getText() + """);

        Comment comment = new Comment(201, userB, "Congratulations!");
        post.setComments(new ArrayList<>());
        post.getComments().add(comment);
        System.out.println("💬 " + userB.getName() + " commented: "" + comment.getText() + """);

        // Auto-notify User A of comment
        Notification commentNotif = new Notification(1004, userB.getName() + " commented on your post.");
        commentNotif.sendNotification(userA);

        // Post summary
        System.out.println("
📢 Post Summary:");
        System.out.println("Post by: " + post.getPostOwner().getName());
        System.out.println("Text: " + post.getText());
        System.out.println("Total Comments: " + post.getComments().size());
        System.out.println("Comment by: " + comment.getCommentOwner().getName());
    }
}
\end{lstlisting}

\subsection{Wrapping up}\label{6jFctsKze0OJS6S6Cfrrq}
\\
We've explored the complete design of LinkedIn in this chapter. We 've looked at how LinkedIn can be visualized using various UML diagrams and designed using object-oriented principles and design patterns.

% End of section content